% LSPV-4 REPORT::PADINV1::2026-10-01T18:49:56-07:00
% !TEX program = xelatex
\documentclass[11pt]{article}
\usepackage[margin=0.88in]{geometry}
\usepackage{amsmath,amssymb}
\begin{document}
\section*{LSPV-4: padding-invariance attack on Search--G counting escapes}
\textbf{STATUS: OPEN.} P versus NP remains open. No unrestricted SAT lower bound is proved here.

\paragraph{Recovery diagnosis and predecessor preflight.}
The repository still has REACQ1 as the newest verified LSPV-4 artifact. The exact LaTeX body and exact connector error of the immediately missing later cycle are not available in the preserved run state accessible here. Per AGENTS.md I classify that immediate recovery cause \textbf{UNKNOWN}; I do not reconstruct or retry missing science. The newest directly checked LSPV-3 report NOGUARD1 is substantive, correctly headed, and has the required verification sentinel, so predecessor preflight does not block this run.

\paragraph{CORPUS REVIEWED.}
I surveyed the current repository directories and governing files, read AGENTS.md, CODEX.md, the pnp-research skill, root main.tex, current task/research directory surfaces, NOGUARD1, REACQ1, RANDCLOSE1, PCLOSE1-R1, and CCSPC1, and reconciled the active SU/vMS/Hadamard, ACQ1, reconstruction, HEC/CRS--PPC/CCSPC/RCCN/Search--G, arbitrary-algorithm, and barrier claims against the living status ledger. Historical recovery failures are not treated as scientific evidence.

\paragraph{Two obstacles and literature preflight.}
The immediate obstacles are (i) whether moving Search--G from one-witness/all-witness predicates to exact counting or parity supplies a representation-stable resource, and (ii) whether such a resource is forced on an arbitrary FP selector rather than merely appended as a harder auxiliary task. Primary literature checked includes Valiant's \#P framework, Toda's PP/PH theorem and Toda--Ogiwara counting-class results, plus Aaronson--Wigderson's algebrization barrier. These establish the significance of counting classes, but do not establish a selector-to-counting forcing theorem for this project.

\paragraph{PROVED/EXACT: selector-equivalent fiber blow-up.}
Let $R(x,y)$ be a polynomially balanced polynomial-time search relation and let $p$ be any polynomial. Define
\[
R_p(x,(y,z))\iff R(x,y),\qquad z\in\{0,1\}^{p(|x|)}.
\]
Then
\[
W_{R_p}(x)=W_R(x)\times\{0,1\}^{p(|x|)},\qquad
|W_{R_p}(x)|=2^{p(|x|)}|W_R(x)|.
\]
Nevertheless $R$ and $R_p$ are selector-equivalent with polynomial overhead. If $S$ selects a witness for $R$, then
\[
S_p(x)=(S(x),0^{p(|x|)})
\]
selects one for $R_p$. Conversely, projection of the first component of any selector for $R_p$ selects a witness for $R$. Thus polynomial search equivalence does not preserve raw witness count.

\paragraph{PROVED/EXACT: parity destruction and normalized-fraction guard.}
Taking $p(n)\ge1$ makes
\[
|W_{R_p}(x)|\equiv0\pmod 2
\]
for every $x$, irrespective of the original parity. Hence raw witness parity is not invariant under this semantics-preserving COPY wrapper. More generally, exact count can be multiplied by an arbitrary power of two while preserving selector complexity.

If an audit $h(x,y)$ ignores the padding coordinate and $h_p(x,(y,z))=h(x,y)$, then both numerator and denominator of the semantic fraction
\[
\frac{|\{y:R(x,y)\wedge h(x,y)\}|}{|W_R(x)|}
\]
are multiplied by the same $2^{p(|x|)}$. That normalized fraction is padding-invariant. This does not make it hard; it only survives this particular representation attack.

\paragraph{Consequence for Search--G.}
The generic chain
\[
P=NP\Rightarrow FP=FNP\Rightarrow \text{selector}\Rightarrow
\text{many witnesses}\Rightarrow\text{count/parity hardness}
\]
is invalid without an additional construction-specific theorem. Valiant's \#P and Toda's counting results do not imply that an arbitrary selector must compute a raw fiber cardinality. A surviving counting route needs both (a) a representation-stable aggregate attached to the same semantic search problem and (b) a selector-to-aggregate forcing theorem with all costs charged. Merely choosing an audit in \#P, PP, or $\oplus P$ imports a stronger problem rather than deriving it from FNP.

\paragraph{Hostile controls and P=NP-compatible behavior.}
COPY is decisive here; DEAD and DELAY can also inflate records without changing semantic search. INTERP/RECOMP can select a canonical representative. Static reread and REACQ1 rerun remain polynomial for deterministic FP selectors. Gaussian/Tseitin, Horn, 2-SAT and bounded-treewidth controls warn that large witness multiplicity or long traces can coexist with efficient decision/search. These behaviors are compatible with hypothetical $P=NP$ and are not evidence for $P=NP$.

\paragraph{PQ/cross-domain classification.}
\textbf{EXACT OBJECT/THEOREM TRANSFER:} fixed-code Schur products, quadratic-relation kernels, typed puncturing/shortening, collision Cauchy--Schwarz, and the padding identities above. \textbf{ANALOGY/CANDIDATE:} syzygies/Betti invariants as succinct separators, special-soundness collision extraction, LWE-style external extraction, and normalized collision entropy as a representation-stable candidate. \textbf{FAILED TRANSFER:} cryptographic hardness/random oracles, VLSI/distributed locality, online no-reread bounds, and proof-system-specific lower bounds as unrestricted static SAT runtime claims.

\paragraph{SU/vMS, acquisition, reconstruction, and barriers retained.}
No claim here changes the source-sanctioned Hadamard branch or exact fixed-code Schur algebra. Exact SU09/vMS invocation tracing and global stacked rank/common kernel remain OPEN. ACQ1 remains CONDITIONAL on the exact legal-state evaluator and model-matched SAT compilation exponent. Sparse-$Q$ normalization and NOGUARD1's quantitative reconstruction improvement remain OPEN. The padding theorem relativizes, and the active Schur/polar machinery remains algebraic; Aaronson--Wigderson's nonalgebrization barrier remains live.

\paragraph{DUPLICATION / CONTRADICTION AUDIT.}
REACQ1 already retired generic superpolynomial reacquisition for deterministic polynomial selectors. PCLOSE1-R1 and RANDCLOSE1 already retired deterministic and bounded-error randomized polynomial-depth causal composition as standalone escapes. The new contribution is orthogonal: it attacks raw counting/parity as representation-dependent properties of an FNP witness relation. No contradiction with RCCN's declared STORE-or-REACQUIRE theorem or NOGUARD1's reconstruction accounting was found.

\paragraph{STALE OR OVERSTATED MAIN CLAIMS; adopted/rejected/retired.}
Root main.tex correctly marks P versus NP OPEN and already retires generic STORE-or-REACQUIRE promotion. Any future synthesis that treats raw witness multiplicity, exact fiber count, or witness parity as an invariant hardness resource of semantic FNP search should be rejected. Adopt the source-sanctioned fixed-Hadamard algebra, charged ACQ1, and NOGUARD1 baseline. Retire raw witness count/parity as generic selector-hardness resources.

\paragraph{Strict status and next target.}
\textbf{PROVED/EXACT:} selector-preserving free-witness padding, multiplicative count distortion, parity destruction, and preservation of padding-independent normalized witness fractions. \textbf{KNOWN THEOREM:} Valiant's \#P framework; Toda and Toda--Ogiwara counting results. \textbf{CANDIDATE:} a canonical representation-invariant aggregate plus a selector-to-aggregate forcing theorem. \textbf{FALSIFIED/RETIRED:} raw witness multiplicity/count/parity as generic FNP selector-hardness resources. \textbf{OPEN:} selector-to-aggregate forcing, exact SU09/vMS invocation, stacked Schur rank, ACQ1 exponent, sparse-$Q$ normalization, reconstruction improvement, arbitrary-algorithm transfer, and the nonrelativizing/nonalgebrizing hinge.

The next falsification target is to apply this padding equivalence to the strongest concrete Search--G aggregate in the repository. If the purported resource changes under $R\mapsto R_p$ while semantic search and selector complexity do not, retire it. Only a padding-invariant aggregate should proceed to a forcing-theorem attempt.

\begin{thebibliography}{4}
\bibitem{Valiant79} L. G. Valiant, \emph{The Complexity of Enumeration and Reliability Problems}, SIAM J. Comput. 8(3), 410--421, 1979.
\bibitem{Toda91} S. Toda, \emph{PP is as Hard as the Polynomial-Time Hierarchy}, SIAM J. Comput. 20(5), 865--877, 1991.
\bibitem{TodaOgiwara92} S. Toda and M. Ogiwara, \emph{Counting Classes are at Least as Hard as the Polynomial-Time Hierarchy}, SIAM J. Comput. 21(2), 316--328, 1992.
\bibitem{AW09} S. Aaronson and A. Wigderson, \emph{Algebrization: A New Barrier in Complexity Theory}, ACM Trans. Comput. Theory 1(1), 2009.
\end{thebibliography}
\end{document}
% END_OF_VERIFIED_REPORT::LSPV-4::PADINV1